\hbadness=9000
% Sprachunterstützung 
% Mehrere Sprachen für angepasstes Literaturverzeichnis (babelbib) 
\usepackage[english,german,ngerman]{babel}
% Internationales Literaturverzeichnis
\usepackage{babelbib}
\setbtxfallbacklanguage{english}




% http://mrunix.de/forums/showthread.php?p=301707
% Abstand zwischen Absätzen in minipage-Umgebung:
\makeatletter
\newcommand*{\@minipagerestore}{\parskip@update}
\makeatother



%%%%%%%%%%%%%%%%%%%%%%%%%%%%%%%%%%%%%%%%%%%%%%%%%%%%%%%%%%%
%%% Verzeichnisse
\usepackage[
% tight,
]{shorttoc}
%
% % When you want to set only the table of contents of a document in two (or
% % more columns), there is one known way1: Add an
% % \addtocontents{toc}{\protect\begin{multicols}{2}}
% % before the \tableofcontents and an
% % \addtocontents{toc}{\protect\end{multicols}}
% % at the end of the document. This way your .toc will start with
% % \begin{multicols}{2} and end with \end{multicols}.
%
% \usepackage[toc]{multitoc}

% increase width for large page numbers in TOC
%% http://www.kronto.org/thesis/tips/format-toc.html
%% Because of the font change, the page number becomes too large for the
%% horizontal space LaTeX reserves for it by default. Without the following
%% redefines to fix it, this would cause the Chapter entry page numbers
%% to extend a few points into the right margin. The horror!

\makeatletter 
\renewcommand*{\@pnumwidth}{3.5em}
\renewcommand{\@tocrmarg}{4.5em}

\makeatother


% /part ohne Seitenzahl im TOC?
% http://mrunix.de/forums/showthread.php?t=61596
\makeatletter
\let\partbackup\l@part
\renewcommand*\l@part[2]{\partbackup{#1}{}}
\makeatother

% \makeatletter
% \let\chapterbackup\l@chapter
% \renewcommand*\l@chapter[2]{\chapterbackup{#1}{}}
% \makeatother


\usepackage{makeidx}
\setcounter{secnumdepth}{3}
% LaTeX-Begleiter S.676
\makeatletter
\renewenvironment{theindex}
   {
   \begin{WidePar}
   ~

   \begin{multicols}{3}[\chapter{\indexname}]
      \setlength\parindent{0pt}\setlength\parskip{0pt}\let\item\@idxitem
   }
   {
   \end{multicols}
   \end{WidePar}
   }
\makeatother

% Hat man den Wunsch, auf eine Fußnote mehrfach zu verweisen, bietet LATEX dafür keine direkte Möglichkeit.
% Man muss hier die folgende Definition, eingerahmt von \makeatletter und \makeatother, in den Vorspann des Dokuments einfügen.
% Im Text muss bei der Fußnote, auf die ein zweites mal verwiesen werden soll, ein Label (hier: \label{fn:wichtig})eingefügt werden.
% Mit \footref{fn:wichtig} verweist man dann ein weiteres mal auf sie.
% http://www2.informatik.hu-berlin.de/~piefel/LaTeX-PS/Archive-2004/V07-footnote.pdf
\makeatletter
\providecommand*{\footref}[1]{% \begingroup 
\unrestored@protected@xdef\@thefnmark{\ref{#1}}% \endgroup 
\@footnotemark}
\makeatother

\usepackage{prettyref}
\newrefformat{sec}{\emph{\ref{#1}}}
% \newrefformat{pkt}{\emph{Abs.\,\ref{#1}}}
% \newrefformat{lit}{\emph{Abs.\,\ref{#1}}}
\newrefformat{chp}{Kap.~\ref{#1}}
\newrefformat{fig}{\SignFigureSmall\,\ref{#1}}
\newrefformat{tab}{\SignTableSmall\,\ref{#1}}
% \newrefformat{tab}{{\sffamily\mdseries\color{DefaultB}{\scalefont{0.7}\ovalbox{T}}}\,\ref{#1}}
% \newrefformat{m}{{\sffamily\mdseries\color{DefaultB}{\scalefont{0.7}\ovalbox{M}}}\,\ref{#1}}
\newrefformat{m}{\SignMassnahmeSmall\,\ref{#1}}
\newrefformat{topic}{\Pointinghand\,\ref{#1}}

\usepackage[AASS,stylelayout]{aass-common}
\usepackage{aass-units}
\usepackage{aass-fonts}
\usepackage{aass-features}
\usepackage{aass-text}










%%%%%%%%%%%%%%%%%%%%%%%%%%%%%%%%%%%%%%%%%%%%%%%%%%%%%%%%%%%
%%% Seitenlayout
% \usepackage[
% verbose,
% a4paper,
% asymmetric,
% % reversemp,
% includemp,   % incl. marginpar
% includehead, % incl. header
% includefoot, % incl. footer
% inner  = 11mm,
% outer  = 11mm,
% top    = 10mm,
% bottom =  8mm,
% bindingoffset=15mm,
% marginparwidth = 60mm,
% footskip = 7mm,
% % % headheight=0.5cm,
% % % headsep=0.5cm,
% % showframe,
% ]
% {geometry}


% \usepackage[
% verbose,
% a4paper,
% asymmetric,
% % reversemp,
% % includemp,
% % tmargin=25mm,
% % bmargin=20mm,
% lmargin=13mm,
% rmargin=70mm,
% % outer  = 56mm,
% % inner  = 56mm,
% top    = 20mm,
% bottom = 15mm,
% % % hmarginratio=2:1,
% marginparwidth = 60mm,
% % % headheight=0.5cm,
% % % headsep=0.5cm,
% footskip = 5mm,
% % includehead,
% bindingoffset=1.5cm,
% showframe
% ]
% {geometry}




% \usepackage{scrindex}
\makeindex






%%% Tables

\usepackage{boxedminipage}

\usepackage{layout}
% \usepackage{calc}
% \usepackage[strict]{changepage}
%\usepackage{gmeometric}
% \eometry{verbose,
% a4paper,
% asymmetric,
% % tmargin=25mm,
% % bmargin=20mm,
% % lmargin=56mm,
% % rmargin=56mm,
% % outer  = 56mm,
% % inner  = 56mm,
% % top    = 25mm,
% % bottom = 20mm,
% marginparwidth = 40mm,
% % headheight=0.5cm,
% % headsep=0.5cm,
% footskip = 10mm,
% % includehead,
% includemp,
% bindingoffset=1.5cm
% }


\newlength{\Textspalte}
\newlength{\Randspalte}
\newlength{\Trenner}
\newlength{\Korrektur}
\setlength{\Textspalte}{\textwidth}
\setlength{\Randspalte}{\marginparwidth}
\setlength{\Trenner}{6pt}
\setlength{\Korrektur}{{\fboxrule} + {\fboxsep}}


%%%%%%%%%%%%%%%%%%%%%%%%%%%%%%%%%%%%%%%%%%%%%%%%%%%%%%%%%%%
%%% Pargragraphensatzregeln
%Dann darst du die parskip und \parindent gar nicht verändern, sondern Koma Script in der Klasse die entsprechende Option übergeben. 
%\setlength{\parindent}{0pt}
%\setlength{\parskip}{5pt}
\raggedbottom
\tolerance=9999
\emergencystretch=20pt
%schusterjungen
% Disable single lines at the start of a paragraph (Schusterjungen)
%\clubpenalty = 10000
%
% Disable single lines at the end of a paragraph (Hurenkinder)
%\widowpenalty = 10000 \displaywidowpenalty = 10000

\usepackage{fancybox}


% User page numbering with chapter number in it

% \renewcommand{\thechapter}{\Alph{chapter}}

% \renewcommand{\thechapter}{{\ttfamily\Alph{chapter}}}
% \renewcommand{\thesection}{{\color{DefaultA}\thechapter .\arabic{section}}}
% \renewcommand{\thesubsection}{{\color{DefaultA}\thesection .\arabic{subsection}}}
% \renewcommand{\thesubsection}{{\color{aassgrey}\thesection.}\arabic{subsection}}
\renewcommand{\thesubsubsection}{{\thesubsection .\alph{subsubsection}}}



% No bottom titles !!!
\usepackage[nobottomtitles]{titlesec}
\renewcommand{\bottomtitlespace}{.10\textheight}

% \titleformat{\section}[rightmargin]{\usefont{T1}{lmss}{bx}{n}\selectfont\LARGE\color{DefaultA}}{}{0em}{}
% \titleformat{\section}{\large\scshape\raggedright}{}{1em}{}


%%%%%%%%%%%%%%%%%%%%%%%%%%%%%%%%%%%%%%%%%%%%%%%%%%%%%%%%%%%
%%% Floats
\usepackage{wrapfig,picins}

% watches margin paragraphs and
% gives warnings
%\usepackage{mparhack} 

% \usepackage[algo2e,algochapter]{algorithm}
% \usepackage{listings}




\usepackage{hyperref}

% \newcommand{\url}[1]{\ttfamily #1}

% floats waren hier

% \usepackage[grey,helvetica]{quotchap}
% %\renewcommand{\sectfont}{\large\color{DefaultA}\sffamily\fontfamily{phv}\selectfont}
% \definecolor{chaptergrey}{rgb}{0.788235294,0,0.08627451}
% \definecolor{chaptergrey}{named}{DefaultContrast}
% % \definecolor{chaptergrey}{rgb}{0.788235294,0,0.08627451}
% % \newenvironment{savequote}{}{}
% 
% % \usepackage[Rejne]{fncychap}
% % \ChNameVar{\centering\LARGE\sf} 
% % \ChNumVar{\Huge\color{DefaultB}} 
% % \ChTitleVar{\centering\Large\sf\color{DefaultA}}
% % \ChRuleWidth{2pt} 
% % \ChNameUpperCase
% % % \ChNumVar{\fontsize{76}{80}\usefont{OT1}{lmdh}{m}{n}\selectfont\color{DefaultB}}
% \addtokomafont{chapter}{\Huge\sffamily\bfseries\color{DefaultA}}

% \addtokomafont{chapter}{\fontfamily{lmdh}\selectfont\Huge\color{DefaultA}}



% % http://texblog.net/latex/layout/
% \usepackage{tikz}
% % \usepackage{kpfonts}
% % \usepackage[explicit]{titlesec}
% \newcommand*\chapterlabel{}
% \titleformat{\chapter}
%   {\def\chapterlabel{}
%    \normalfont\sffamily\Huge\bfseries\scshape}
%   {\def\chapterlabel{\thechapter\ lala}}{0pt}
%   {\begin{tikzpicture}[remember picture,overlay]
%     \node[yshift=-3cm] at (current page.north west)
%       {\begin{tikzpicture}[remember picture, overlay]
%         \draw[fill=LightSkyBlue] (0,0) rectangle
%           (\paperwidth,3cm);
%         \node[anchor=east,xshift=.9\paperwidth,rectangle,
%               rounded corners=10pt,inner sep=11pt,
%               fill=MidnightBlue]
%               {\color{white}\chapterlabel#1};
%        \end{tikzpicture}
%       };
%    \end{tikzpicture}
%   }

% \titlespacing*{\chapter}{0pt}{50pt}{-60pt}
% 
% \usepackage[automark]{scrpage2}
% \usepackage{tikz}
% \newcommand*\tikzhead[1]{%
%   \begin{tikzpicture}[remember picture,overlay]
%     \node[yshift=-2cm] at (current page.north west)
%       {\begin{tikzpicture}[remember picture, overlay]
%         \draw[fill=LightSkyBlue] (0,0) rectangle
%           (\paperwidth,2cm);
%         \node[anchor=east,xshift=.9\paperwidth,rectangle,
%               rounded corners=15pt,inner sep=11pt,
%               fill=MidnightBlue]
%               {\color{white}#1};
%        \end{tikzpicture}
%       };
%    \end{tikzpicture}}
% \clearscrheadings
% \ihead{\tikzhead{\headmark}}
% \pagestyle{scrheadings}
% 
% 
% \usepackage[automark]{scrpage2}
% 
% \usepackage{tikz}
% \newcommand*\tikzhead[1]{%
%   \begin{tikzpicture}[remember picture,overlay]
%     \node[yshift=-2cm] at (current page.north west)
%       {\begin{tikzpicture}[remember picture, overlay]
%         \draw[fill=LightSkyBlue] (0,0) rectangle
%           (\paperwidth,2cm);
%         \node[anchor=east,xshift=.9\paperwidth,rectangle,
%               rounded corners=15pt,inner sep=11pt,
%               fill=MidnightBlue]
%               {\color{white}#1};
%        \end{tikzpicture}
%       };
%    \end{tikzpicture}}
% \clearscrheadings
% \ihead{\tikzhead{\headmark}}
% \pagestyle{scrheadings}
% 
% 


% \newenvironment{savequote}{}{}







\usepackage{minitoc} 
%The quotchap package should be loaded BEFORE the minitoc package.
\newcounter{DefaultMinitocDepth}
\setcounter{DefaultMinitocDepth}{3}
% \mtcsetdepth{minitoc}{\value{DefaultMinitocDepth}}

\nomtcrule

% \setlength{\mtcindent}{0pt}
% \setlength{\stcindent}{0pt}
% \mtcindent{0pt}
% \stcindent{0pt}
% \renewcommand{\mtcfont}{\sffamily\small}
% \renewcommand{\mtcSfont}{\sffamily\small\bfseries}
% \renewcommand{\mtcSSfont}{\sffamily\scriptsize}
% \renewcommand{\mtcSSSfont}{\sffamily\scriptsize}
% \mtcsetoffset{minitoc}{-4em}
% \mtcsetformat{mini-table}{mtcindent}{0pt}

\setlength{\mtcindent}{0pt}
\mtcsetfeature{minitoc}{open}{\kern1sp\vspace*{-.1ex}\begin{multicols}{2}[\kern-2.5ex]}
\mtcsetfeature{minitoc}{close}{\end{multicols}\kern-2.ex}
% \mtcsetfeature{minitoc}{open}{\begin{multicols}{2}}
% \mtcsetfeature{minitoc}{close}{\end{multicols}}
\mtcsetfeature{minitoc}{before}{\begin{WidePar}\raggedcolumns}
\mtcsetfeature{minitoc}{after}{\flushcolumns\end{WidePar}}
\mtcsetoffset{minitoc}{-0.7em}
\setlength{\mtcindent}{-0.7em}
% \mtcsetformat{minitoc}{tocrightmargin}{2cm}

% % multitoc is incompatible with float package
% \usepackage[toc]{multitoc}
% \renewcommand{\multicolumntoc}{2}







% \usepackage[noauto]{chappg} % nach hyperref!
% 
% \pagenumbering{bychapter}
% \renewcommand{\chappgsep}{/}


%%%%%%%%%%%%%%%%%%%%%%%%%%%%%%%%%%%%%%%%%%%%%%%%%%%%%%%%%%%
%%% Beschriftungen
\DeclareRobustCommand{\SignTable}{{\noindent\sffamily\mdseries\color{DefaultB}\ovalbox{T}}}
\DeclareRobustCommand{\SignTableSmall}{{\scalefont{0.7}\SignTable}}

\DeclareRobustCommand{\SignFazit}{\noindent{\color{DefaultContrastA}\bfseries{\VarFlag ~}}}
\DeclareRobustCommand{\SignFazitSmall}{\scalefont{0.7}\SignFazit}

\DeclareRobustCommand{\SignLeitsatz}{{\noindent\sffamily\mdseries\color{DefaultB}\ovalbox{L}}}
\DeclareRobustCommand{\SignLeitsatzSmall}{{\scalefont{0.7}\SignLeitsatz}}

\DeclareRobustCommand{\SignFigure}{{\noindent\sffamily\mdseries\color{DefaultB}\ovalbox{A}}}
\DeclareRobustCommand{\SignFigureSmall}{{\scalefont{0.7}\SignFigure}}

\DeclareRobustCommand{\SignMassnahme}{{\noindent\sffamily\mdseries\color{DefaultB}\ovalbox{M}}}
\DeclareRobustCommand{\SignMassnahmeSmall}{{\scalefont{0.7}\SignMassnahme}}









%%%%%%%%%%%%%%%%%%%%%%%%%%%%%%%%%%%%%%%%%%%%%%%%%%%%%%%%%%%
%%% Bibliographie

% Bibliographie-Stile babxxx unterstützen internationale Bibliographien, Paket
% 'babelbib'.
% Die Sprache muss in BibTeX-Datei im Feld 'language' stehen (dzt. erlaubt:
% english, german, ngerman). Die entsprechende Sprache muss in babel geladen
% worden sein.
\bibliographystyle{babunsrt-fl}
% \bibliographystyle{bababbrv-fl}
% \bibliographystyle{plainnat}
% \usepackage{cite}
\usepackage{bibentry}
% \nobibliography*
\usepackage[
% round, % (default) for round parentheses;
% square, % for square brackets;
% curly, % for curly braces;
% angle, % for angle brackets;
% colon, % (default) to separate multiple citations with colons;
% comma, % to use commas as separaters;
% authoryear, % (default) for author-year citations;
numbers, % for numerical citations;
% super, % for superscripted numerical citations, as in Nature;
% sort, % orders multiple citations into the sequence in which they appear in the
% % list of references;
sort&compress, % as sort but in addition multiple numerical citations are compressed if possible ;
% longnamesfirst, % makes the first citation of any reference the equivalent of
% % the starred variant (full author list) and subsequent citations normal (abbreviated list);
% sectionbib, % redefines \thebibliography to issue \section* instead of
% % \chapter*; valid only for classes with a \chapter command; to be used with the chapterbib package;
% nonamebreak, % keeps all the authors' names in a citation on one line; causes
% % overfull hboxes but helps with some hyperref problems.
]{natbib}
\renewcommand{\citenumfont}[1]{\sffamily\small#1}
% \renewcommand{\citenumfmt}[1]{\sffamily\textbf{#1}:}
\renewcommand{\bibnumfmt}[1]{\sffamily\textbf{#1}:}


% \setcitestyle{
% super,
% open={\sffamily\tiny[},
% close={\sffamily\tiny]},
% comma,
% % aysep={~},
% notesep={text}
% }

% \usepackage[sectionbib]{chapterbib} % Literaturverzeichnisse für Kapitel
% 
% % \ChapterBibliography muss in der jeweiligen Kapiteldatei aufgerufen werden
% % um dort ein entsprechend formatiertes Literaturverzeichnis zu erstellen.
% \newcommand{\ChapterBibliography}
% 	{\bibliographystyle{babunsrt-fl}
% 	    {
%     \loadgeometry{FullPage}\twocolumn
% %     \sffamily\itshape
% %       \begin{multicols}{2}
% 
%         \scriptsize
%         \setlength{\parskip}{1pt}
%         \bibliography{Literatur-AASS}
% %       \end{multicols}
% 
% 
%     \restoregeometry
%     }
% %    \end{WidePar}
%    \onecolumn
% % 	\begin{multicols}{2}
% %     	{
% %     		\scriptsize
% %     		\setlength{\parskip}{1pt}
% %     		\bibliography{./Literatur-AASS}
% %     	}
% %     \end{multicols}
%  	}

\newcommand{\ChapterBibliography}{\relax}


\usepackage[version=3]{mhchem}
\mhchemoptions{textfontcommand=\sffamily}
\mhchemoptions{mathfontcommand=\mathsf}









% Test PLAYGROUND

\newcommand{\FullRef}[1]{{\ttfamily\ref{#1} (S.\,\pageref{#1})}}
\newcommand{\FullPrettyRef}[1]{{\prettyref{#1}, Seite \pageref{#1}}}

\newcommand{\TabItemI}[1]{\begin{inparaitem}\item #1 \end{inparaitem}}
\newcommand{\TabItemII}[1]{\begin{inparaitem}\begin{inparaitem}\item#1\end{inparaitem}\end{inparaitem}}
\newcommand{\TabItemIII}[1]{\begin{inparaitem}\begin{inparaitem}\begin{inparaitem}\item#1\end{inparaitem}\end{inparaitem}\end{inparaitem}}

\setlength\columnsep{3em}
\setlength\columnseprule{0.4pt}
\setlength\premulticols{10\baselineskip}
\renewcommand{\columnseprulecolor}{\color{DefaultGrey}}

\setkomafont{footnote}{\sffamily}
\deffootnotemark{\textsuperscript{\sffamily\thefootnotemark}}
% \setlength{\footnotesep}{1pt}
\usepackage[hang,multiple]{footmisc}
\renewcommand\footnotelayout{\sffamily}






